\title{FlashAttention-3: Fast and Accurate Attention with Asynchrony and Low-precision}

\begin{abstract}
Attention forms a critical component within the Transformer architecture, yet it presents a computational bottleneck for scaling large language models and processing extended contexts. The \textsc{FlashAttention} method introduced optimizations to accelerate attention mechanisms on GPUs by significantly reducing memory access operations. Nevertheless, prior approaches—including \textsc{FlashAttention-2}, which achieves only 35\% compute utilization on the H100 GPU—do not fully leverage advanced features available in state-of-the-art hardware. To address this gap, we propose three key strategies for accelerating attention on Hopper GPUs: harnessing the asynchronicity between Tensor Cores and TMA to (1) concurrently execute computation and memory transfers using warp specialization, (2) intertwine block-wise matrix multiplication with softmax computations, and (3) apply block quantization coupled with incoherent processing that exploits the native FP8 precision support. Our proposed solution, \textsc{FlashAttention-3}, delivers 1.5-2.0$\times$ speed improvements on H100 GPUs, reaching throughput up to 740 TFLOPs/s with FP16 (realizing 75\% hardware utilization), and nearly 1.2 PFLOPs/s with FP8. Empirically, we show that FP8 \textsc{FlashAttention-3} attains a 2.6$\times$ reduction in numerical error relative to a standard FP8 attention baseline.
\end{abstract}

\section{Introduction}
\label{sec:intro}

Within the Transformer architecture~\citep{vaswani2017attention}, the central computational constraint lies in the attention mechanism, as calculating self-attention scores between queries and keys incurs a quadratic complexity relative to the sequence length. Enhancing the scalability of attention mechanisms to handle longer contexts has the potential to facilitate advanced capabilities, such as modeling and reasoning across multiple lengthy documents~\citep{guo2021longt5,shaham2022scrolls,peng2023yarn} or processing extensive code repositories~\citep{roziere2023code, li2023starcoder}, as well as supporting additional data modalities including high-resolution images~\citep{chen2022scaling}, audio~\citep{gulati2020conformer}, and video~\citep{ho2022video}. Furthermore, it enables applications involving prolonged user histories~\citep{sun2019bert4rec} or complex agent workflows over extended timeframes~\citep{yao2022react}. Consequently, there is substantial interest in accelerating attention for long-context scenarios through various means, such as approximate methods~\citep{katharopoulos2020transformers,choromanski2020rethinking,
tay2020efficient}, software-level improvements (\citep{rabe2021self,dao2022flashattention,kwon2023efficient}), or even the exploration of alternative model architectures~\citep{peng2023rwkv,sun2023retentive,gu2023mamba}.

This study expands upon the methodologies proposed by \citet{dao2022flashattention}, which focus on the development of precise attention mechanisms by explicitly incorporating GPU execution paradigms and architectural constraints into algorithmic design. Dao et al., in their work \citep{dao2022flashattention}, introduced \textsc{FlashAttention}, which leverages a unique tile-based approach for efficiently parallelizing attention computation. By consolidating all attention-related procedures into a single GPU kernel, \textsc{FlashAttention} significantly reduces the need for slow global memory accesses, avoiding unnecessary intermediate storage operations.

Subsequently, \citet{dao2023flashattention2} revised this approach with \textsc{FlashAttention-2}, enabling parallelization not only across tiles but also along the sequence length. This updated algorithm processes the inner loop of the forward phase using block-wise computation on the key and value matrices, leading to improved GPU workload distribution and occupancy. Despite these enhancements, our analysis indicates that \textsc{FlashAttention-2} exhibits suboptimal performance utilization on the most recent GPU hardware, especially when compared to finely-tuned matrix multiplication (GEMM) kernels; for instance, only achieving about 35\% utilization on Hopper H100, compared to 80-90\% attained by optimized GEMM implementations. Such discrepancies may be partly explained by implementation specifics, notably the use of instructions intended for previous generation (Ampere) Tensor Cores rather than those tailored for Hopper architecture. Prior works, including ThunkerKitten~\citep{spector2024thunder} and cuDNN 9~\citep{cudnn9}, have demonstrated that adoption of Hopper-specific instructions and tiling abstractions not only accelerates attention operations but also leads to simplification in code structure.

At a foundational level, the algorithm underlying \textsc{FlashAttention-2} is structured around a basic synchronous computation paradigm, without incorporating asynchronous processing or leveraging low-precision arithmetic in its architecture. Asynchrony arises due to specialized hardware components intended to boost performance in key machine learning operations—dedicated units such as Tensor Cores for matrix multiplication or Tensor Memory Accelerators (TMA) for data movement operate independently from CUDA cores, which handle miscellaneous logic as well as integer and floating-point operations. The employment of reduced numerical precision formats—including FP8 on Hopper and FP4 on Blackwell, as successors to the adoption of FP16 (Pascal, 2017) and BF16 (Ampere, 2020)—demonstrates a validated strategy to achieve two or four times greater throughput within the same constraints of power and silicon area. In \cref{subsec:hardware}, we examine Hopper's provisions in these areas. The central technical hurdle involves refactoring \textsc{FlashAttention-2} to exploit these hardware advancements: effective use of asynchrony means aligning matmul and softmax computations concurrently despite dependency relationships, while utilizing low-precision demands strategic handling to prevent excessive quantization error, particularly for rare or extreme-value features in LLMs~\citep{dettmers2208llm, sun2024massive}.

In pursuit of this goal, we introduce \textsc{FlashAttention-3}, which integrates three innovative concepts to enhance efficiency on advanced GPU platforms:\footnote{Although our evaluation centers on NVIDIA's Hopper architecture, the algorithm remains applicable to any GPU architecture that possesses strong support for asynchronous operations and low-precision computation.}

\begin{enumerate}

\item \textbf{Producer-Consumer asynchrony:} We implement a warp-specific software pipelining approach that takes advantage of asynchronous data transfers and Tensor Core operations by assigning data producers and consumers to different warps. This partitioning effectively broadens the algorithm's capacity to mask latencies associated with both memory access and instruction issuance.

\item \textbf{Hiding softmax under asynchronous block-wise GEMMs:} To increase throughput, we orchestrate the non-GEMM steps in softmax, such as floating-point fused multiply-add and exponentiation, to execute in tandem with the asynchronous WGMMA GEMM instructions. In doing so, we modify the \textsc{FlashAttention-2} algorithm to avoid sequential dependencies between softmax and the GEMM phases. For instance, in our 2-stage variant, softmax computation proceeds on one block of the score matrix while WGMMA, via asynchronous proxy, processes the subsequent block.

\item \textbf{Hardware-accelerated low-precision GEMM:} We reconfigure the forward pass to utilize FP8 Tensor Cores for GEMM operations, achieving close to twice the observed TFLOPs/s. This adaptation demands managing WGMMA’s particular memory layout constraints, specifically regarding FP32 accumulators and FP8 operands. To counteract potential precision degradation from using FP8, we apply block quantization techniques and incoherent processing.

To assess our approach experimentally, we evaluate \textsc{FlashAttention-3} on the H100 SXM5 GPU across various settings. Our results indicate that: (1) FP16 attains a speedup ranging from 1.5$\times$ to 2.0$\times$ in the forward pass compared to \textsc{FlashAttention-2}, peaking at 740 TFLOPs/s, and exhibits a 1.5-1.75$\times$ acceleration in the backward pass; (2) FP8 reaches nearly 1.2 PFLOPs/s; and (3) at longer sequence lengths, FP16 surpasses alternatives while FP8 offers competitive performance\footnote{Specifically, FP8 is superior for head dimension 64 in \textsc{FlashAttention-3}, is on par with higher dimensions (128 and 256) absent causal masking, and trails when causal masking is applied.} relative to the latest attention module in NVIDIA’s cuDNN library. Furthermore, FP16 in \textsc{FlashAttention-3} maintains equivalent numerical accuracy to \textsc{FlashAttention-2} and outperforms the baseline attention mechanism, since intermediate computations (such as softmax normalization) remain in FP32. Additionally, employing block quantization and incoherent execution in FP8 with \textsc{FlashAttention-3} yields a 2.6$\times$ accuracy improvement compared to traditional attention using per-tensor quantization, especially when encountering outlier feature values.

We have released \textsc{FlashAttention-3} under a permissive license\footnote{\textsc{FlashAttention-3}
  is available at \texttt{https://github.com/Dao-AILab/flash-attention}} and intend to incorporate it into both PyTorch and Hugging Face libraries, aiming to maximize its accessibility and utility for researchers and practitioners.

\section{Background: Multi-Head Attention and GPU Characteristics}
\label{sec:background}

\subsection{Multi-Head Attention}
\label{subsec:multi_head_attn}

Let $\mathbf{Q}, \mathbf{K}, \mathbf{V} \in \mathbb{R}^{N \times d}$ be the query, key and value input sequences associated to a single head, where $N$ is the sequence length and $d$ is the head dimension. Then the attention output $\mathbf{O}$ is computed as:

\begin{equation*}
  \mathbf{S} = \alpha \mathbf{Q} \mathbf{K}^\top \in \mathbb{R}^{N \times N}, \quad \mathbf{P} = \mathrm{softmax}(\mathbf{S}) \in \mathbb{R}^{N \times N}, \quad \mathbf{O} = \mathbf{P}\mathbf{V} \in \mathbb{R}^{N \times d},
\end{equation*}

Here, $\mathrm{softmax}$ operates on each row, and conventionally, the scaling parameter $\alpha$ is chosen as $1/\sqrt{d}$. To mitigate numerical issues associated with the exponential computation, it is common to subtract $\mathrm{rowmax}(\mathbf{S})$ from $\mathbf{S}$. In the context of multi-head attention (MHA), separate query, key, and value transformations are allocated for each head. The attention computation is carried out simultaneously over various heads and batches, resulting in the aggregate output tensor.

Now let $\phi$ be a scalar loss function and let $\mathbf{d}(-) = \partial \phi / \partial (-)$ be notation for the gradient. Given the output gradient $\mathbf{dO} \in \mathbb{R}^{N \times d}$, we compute $\mathbf{dQ}$, $\mathbf{dK}$, and $\mathbf{dV}$ according to the chain rule as follows:

\begin{align*}
  \mathbf{dV} &= \mathbf{P}^\top \mathbf{dO} \in \mathbb{R}^{N \times d} \\
  \mathbf{dP} &= \mathbf{dO} \mathbf{V}^\top \in \mathbb{R}^{N \times N} \\
  \mathbf{dS} &= \mathrm{dsoftmax} (\mathbf{dP}) \in \mathbb{R}^{N \times N} \\
  \mathbf{dQ} &= \alpha \mathbf{dS} \mathbf{K} \in \mathbb{R}^{N \times d} \\
  \mathbf{dK} &= \alpha \mathbf{dS}^\top \mathbf{Q} \in \mathbb{R}^{N \times d},
\end{align*}

In this context, we observe that for $p = \mathrm{softmax}(s)$ with respect to a vector $s$, the following relationship holds: $\mathbf{d}s = (\mathrm{diag}(p) - p p^\top)\mathbf{d}p$. We denote by $\mathrm{dsoftmax}(\mathbf{dP})$ the row-wise application of this expression. Ultimately, this calculation is also amenable to parallelization over multiple heads and batch elements during the backward pass in MHA.

\subsection{GPU hardware characteristics and execution model}
\label{subsec:hardware}

We outline the elements of the GPU execution model that pertain to \textsc{FlashAttention-3}, concentrating specifically on the NVIDIA Hopper architecture as an illustrative example of this framework.

\paragraph{Memory hierarchy:} GPU memory is structured into multiple layers of data storage, where higher capacity comes at the cost of reduced bandwidth (\cref{tab:gpu-hierarchy})\footnote{\citet{luo2024benchmarking} indicates shared memory bandwidth is 128 bytes per clock cycle per SM; by aggregating over 132 SMs and considering a boost clock of 1830 MHz, we estimate total bandwidth.}. The lowest level, Global memory (GMEM)—or HBM—consists of off-chip DRAM accessible by all streaming multiprocessors (SMs). Any data fetched from GMEM is automatically cached in the on-chip L2 cache. Beyond this, each SM has a small, highly banked, and explicitly managed cache called shared memory (SMEM). The innermost tier consists of the register file located within each SM.

\paragraph{Thread hierarchy:} In the GPU programming paradigm, execution units are logically arranged into threads. This hierarchy ranges from the most granular to the broadest: individual threads, warps (each containing 32 threads), warpgroups (composed of 4 adjacent warps), threadblocks (also known as cooperative thread arrays or CTAs), threadblock clusters (available on Hopper), and finally, grids.

There is a strong interdependence between these two hierarchies. All threads belonging to a given CTA are simultaneously assigned to the same SM, and multiple CTAs within a cluster are concurrently executed on a single GPC. Shared memory (SMEM) is accessible to any thread inside the CTA, while each individual thread is restricted to a maximum of 256 dedicated registers (RMEM) for private use.

\paragraph{Asynchrony and warp-specialization:}

GPUs achieve high throughput by leveraging parallelism and asynchronous operations to mask memory and processing delays. In Hopper, the process of asynchronous memory transfer between GMEM and SMEM is facilitated by a specialized hardware component known as the Tensor Memory Accelerator (TMA) \cite[\S7.29]{cuda}. Additionally, Hopper’s Tensor Core, which is accessible through the warpgroup-wide WGMMA instruction \cite[\S9.7.14]{ptx}, departs from earlier architectures like Ampere by operating asynchronously and supporting direct input sourcing from shared memory.

Hardware-assisted asynchrony facilitates the creation of warp-specialized kernels, where the warps within a CTA are explicitly assigned as either producers or consumers, thus executing solely data movement or computational tasks, respectively. This specialization generally enhances the compiler’s capacity to construct superior instruction schedules \citep{warp-specialization-2011}. Furthermore, Hopper introduces the ability to dynamically redistribute registers among warpgroups using \verb|setmaxnreg| \cite[\S9.7.17.1]{ptx}, thereby enabling warps performing MMAs to access a larger portion of RMEM relative to those restricted to TMA operations (which require only one thread per warp).

\paragraph{Low-precision number formats:}
\label{sec:low-precision-gpu}
Recent GPUs are equipped with dedicated hardware blocks designed to optimize operations on low-precision data types. Specifically, the WGMMA instruction is capable of utilizing the FP8 Tensor Cores found in Hopper architecture, effectively doubling the computational throughput per SM relative to the use of FP16 or BF16.

Properly utilizing FP8 WGMMA requires awareness of the operand layout restrictions. In the context of a GEMM operation computing $A \times B^{\top}$, where $A$ has dimensions $M \times K$ and $B$ has dimensions $N \times K$, an operand is considered \emph{mn-major} if its memory arrangement is contiguous along the $M$ (for $A$) or $N$ (for $B$) axis, whereas it is classified as \emph{k-major} if the contiguous memory dimension corresponds to $K$. For FP16 WGMMA, both mn-major and k-major layouts are permitted for operands stored in SMEM. In contrast, FP8 WGMMA restricts support solely to k-major operand organization. Furthermore, for use cases such as attention, where it is desirable to combine consecutive GEMM operations within a unified kernel, incompatibilities between the FP32 accumulator layout and FP8 operand requirements can hinder the sequential use of dependent FP8 WGMMAs.

Within the realm of attention mechanisms, these limitations on layout necessitate adjustments to how an FP8 algorithm is constructed. The specific adaptations are outlined in \cref{sec:algofp8}.

\subsection{Standard Attention and Flash Attention} In accordance with \citet{dao2022flashattention}, the term \textbf{standard attention} refers to an approach for computing attention on the GPU wherein the intermediate matrices $\mathbf{S}$ and $\mathbf{P}$ are explicitly stored in HBM. The primary contribution of \textsc{FlashAttention} was to introduce a local softmax reduction strategy, thereby eliminating the need for costly intermediate memory operations and consolidating the entire attention computation within a single kernel invocation. This localized version of the softmax aligns with lines \ref{code-ws:softmax_start}-\ref{code-ws:softmax_end} of the consumer mainloop presented in \cref{alg:flash3_wgmma_ws_only}, in conjunction with the rescales applied to blocks of $\mathbf{O}$. A concise proof that this method indeed yields $\mathbf{O}$ is provided in \cite[\S 2.3.1]{dao2023flashattention2}.

\section{FlashAttention-3: Algorithm}
\label{sec:algo}

This section outlines the \textsc{FlashAttention-3} algorithm. To streamline exposition, we center our discussion on the forward pass, while the approach for the backward pass is detailed in~\cref{sec:algo_ws_bwd}. Our first step is to demonstrate the incorporation of warp-specialization paired with a circular SMEM buffer into the fundamental structure employed by \textsc{FlashAttention-2}. Next, we detail how WGMMA’s asynchronous behavior can be harnessed to construct a two-stage pipeline, overlapping GEMM and softmax computations. Lastly, we discuss requisite adaptations for FP8, covering layout requirements and enhanced precision achieved through block quantization and incoherent computation.

\subsection{Producer-Consumer asynchrony through warp-specialization and pingpong scheduling}
\label{sec:algo_ws}

\paragraph{Warp-specialization}  
In alignment with \textsc{FlashAttention-2}, the forward computation in \textsc{FlashAttention-3} is highly parallelizable across the batch dimension, head index, and the length of the query sequence. Therefore, a CTA-level abstraction suffices to characterize the algorithm’s operation: it processes a tile $\mathbf{Q}_i$ from the query matrix to yield the corresponding output tile $\mathbf{O}_i$. To streamline our exposition, we initially present the warp-specialization approach employing a circular SMEM buffer, deliberately excluding GEMM-softmax overlap. Let $d$ represent the dimensionality of each head and $N$ the total sequence length, and select a query tile size $B_r$ so that $\mathbf{Q}$ is partitioned into $T_r = \lceil \frac{N}{B_r} \rceil$ blocks, denoted as $\mathbf{Q}_1, .., \mathbf{Q}_{T_r}$.

\begin{algorithm}[H]
    \caption{\small\label{alg:flash3_wgmma_ws_only}\textsc{FlashAttention-3} forward pass \textbf{excluding} intra-consumer overlap -- CTA view}
    \begin{algorithmic}[1]
\REQUIRE Matrices $\mathbf{Q}_i \in \mathbb{R}^{B_r \times d}$, $\mathbf{K}, \mathbf{V} \in \mathbb{R}^{N \times d}$ are stored in HBM, with key block size $B_c$, and total block count $T_c = \lceil \frac{N}{B_c} \rceil$.
\STATE Create a pipeline object to coordinate synchronization barriers using a circular SMEM buffer spanning $s$ stages.
\IF {executing within producer warpgroup}
\STATE Free the specified number of registers in advance.
\STATE Load $\mathbf{Q}_i$ from HBM into shared memory.
\STATE Once loaded, signal to consumers that $\mathbf{Q}_i$ has been transferred.
\FOR{$j = 0$ \textbf{ to } $T_c - 1$}
    \STATE Await consumption of the $(j\,\%\,s)$th buffer stage.
    \STATE Initiate transfer of $\mathbf{K}_j, \mathbf{V}_j$ from HBM to shared memory into the $(j\,\%\,s)$th buffer slot.
    \STATE Upon completion, indicate to consumers that $\mathbf{K}_j$ and $\mathbf{V}_j$ have been loaded.
\ENDFOR
\ELSE 
\STATE Allocate the appropriate number of registers based on the number of consumer warps.
\STATE On-chip, set $\mathbf{O}_i = (0) \in \mathbb{R}^{B_r \times d}$; initialize $\ell_i, m_i = (0), (-\infty) \in \mathbb{R}^{B_r}$.
\STATE Wait until $\mathbf{Q}_i$ becomes available in shared memory.
\FOR{$j = 0$ \textbf{ to } $T_c - 1$}
\STATE Pause until $\mathbf{K}_j$ is available in shared memory.
\STATE Evaluate $\mathbf{S}_i^{(j)} = \mathbf{Q}_i \mathbf{K}_j^T$ using SS-GEMM, then synchronize as necessary. \label{alg:ws_only_gemm1}
\STATE Temporarily save $m_i^{\mathrm{old}} = m_i$, and update $m_i = \mathrm{max}(m_i^{\mathrm{old}}, \mathrm{rowmax}(\mathbf{S}_i^{(j)}))$. \label{code-ws:softmax_start}
\STATE Calculate $\widetilde{\mathbf{P}}_i^{(j)} = \mathrm{exp}(\mathbf{S}_i^{(j)} - m_i)$ and update $\ell_i = \mathrm{exp}(m_i^{\mathrm{old}} - m_i) \ell_i + \mathrm{rowsum}(\widetilde{\mathbf{P}}_i^{(j)})$. \label{code-ws:softmax_end}
\STATE Await readiness of $\mathbf{V}_j$ in shared memory.
\STATE Update $\mathbf{O}_i = \mathrm{diag}(\mathrm{exp}(m_i^{\mathrm{old}} - m_i))^{-1} \mathbf{O}_i + \widetilde{\mathbf{P}}_i^{(j)} \mathbf{V}_j$ via RS-GEMM, then synchronize. \label{alg:ws_only_gemm2}
\STATE Mark the $(j\,\%\,s)$th buffer slot as available for producer use.
\ENDFOR
\STATE Compute final outputs $\mathbf{O}_i = \mathrm{diag}(\ell_i)^{-1} \mathbf{O}_i$ and $L_i = m_i + \log(\ell_i)$.
\STATE Output $\mathbf{O}_i$ and $L_i$ to HBM as the $i$th segments of matrices $\mathbf{O}$ and $L$.
\ENDIF
\end{algorithmic}
\end{algorithm}

In our implementation of \cref{alg:flash3_wgmma_ws_only} targeting Hopper, we employ \verb|setmaxnreg| to manage (de)allocations, utilize TMA for fetching $\mathbf{Q}_i$ as well as the collection $\{ \mathbf{K}_j, \mathbf{V}_j \}_{0 \leq j < T_c}$, and invoke WGMMA to perform GEMM operations within the consumer mainloop—applying the SS or RS prefix depending on whether the first operand resides in shared memory or the register file. When analyzing the control flow in \cref{alg:flash3_wgmma_ws_only}, it is important to recognize that TMA load requests progress asynchronously and thus do not block due to pending loads. Additionally, during the initial $s$ passes of the producer mainloop, wait instructions are absent, allowing the buffer to populate without pauses.

\paragraph{Pingpong scheduling} Due to the asynchronous execution enabled by WGMMA and TMA, combined with warp specialization, it becomes feasible to concurrently process softmax computations for one warpgroup while another warpgroup performs GEMM operations. This approach is warranted by observing that, on current hardware accelerators, non-matmul tasks—such as the computation of exponentials—are significantly less efficient than matmul tasks. For instance, the H100 SXM5 GPU achieves 989 TFLOPS for FP16 matmul operations but only reaches 3.9 TFLOPS for special functions like exponential\footnote{According to the CUDA programming guide, each streaming multiprocessor (SM) is capable of executing 16 special function operations per clock cycle. By multiplying 16 by the total 132 SMs and the clock frequency of 1830 MHz, the result is 3.9 TFLOPS for special functions.}, which are required for softmax. When running an attention forward pass in FP16 with a head dimension of 128, the demand for matmul FLOPS is 512 times greater than that for exponentials, yet the hardware supports exponentials at only a 256-fold lower rate. Consequently, exponentials may consume around 50\% of the cycles used for matmul. The disparity is amplified in FP8 mode: while matmul throughput doubles, exponential throughput remains constant.

Because the exponential operation is handled by a distinct hardware component, the multi-function unit, it is optimal to initiate the exponential computation simultaneously with the matmul execution on the Tensor Cores. To facilitate this overlap, we leverage synchronization barriers (\texttt{bar.sync} instructions), enforcing an execution order where the GEMMs (specifically, GEMM1 -- $\mathbf{P} \mathbf{V}$ for one iteration, and GEMM0 -- $\mathbf{Q} \mathbf{K}^\top$ for the subsequent iteration) associated with warpgroup 1 are completed prior to the GEMMs in warpgroup 2. Consequently, this arrangement enables the softmax computation in warpgroup 1 to run in parallel with GEMMs in warpgroup 2. Subsequently, the two warpgroups alternate, so that warpgroup 2 processes softmax while warpgroup 1 switches back to GEMMs—a mechanism commonly referred to as ``pingpong'' scheduling. This approach is depicted in~\cref{fig:pingpong_scheduling}. While the real-world implementation does not always achieve the perfectly staggered execution shown in the illustration, this method generally boosts performance (for instance, increasing FP16 forward throughput from 570 TFLOPS up to 620–640 TFLOPS for a head dimension of 128 and sequence length of 8192).

\paragraph{Attention variants} In the cases of multi-query attention~\citep{shazeer2019fast} and grouped query attention~\citep{ainslie2023gqa}, we adopt the indexing modifications utilized in \textsc{FlashAttention-2}, thereby preventing unnecessary replication of $\mathbf{K}$ and $\mathbf{V}$ within HBM.

\subsection{Intra-warpgroup overlapping GEMMs and softmax}

Within a single warpgroup, it is possible to execute certain instructions from the softmax concurrently with instructions from the GEMMs. Here, we outline a specific method to achieve this overlap.

In the attention algorithm, operations within the inner loop (main loop) have sequential dependencies that impede parallelization within a single iteration. For example, (local) softmax (lines \ref{code-ws:softmax_start} to \ref{code-ws:softmax_end}) relies on the output $\mathbf{S}_i^{(j)}$ of the first GEMM, while the second GEMM takes its result $\widetilde{\mathbf{P}}_i^{(j)}$ as an operand. Indeed, the wait statements in lines \ref{alg:ws_only_gemm1} and \ref{alg:ws_only_gemm2} of \cref{alg:flash3_wgmma_ws_only} serialize the execution of softmax and GEMMs. However, we can break these dependencies by pipelining across iterations through additional buffers in registers. Pursuing this idea, we propose the following two-stage\footnote{Note that the number of stages of the overlapping scheme is bounded by, but need not equal, the number $s$ of stages in the circular SMEM buffer.} GEMM-softmax pipelining algorithm:

\begin{algorithm}[H]
    \caption{\small\label{alg:flash3_wgmma}\textsc{FlashAttention-3} consumer warpgroup forward pass}
    \begin{algorithmic}[1]
      \REQUIRE  Matrices $\mathbf{Q}_i \in \mathbb{R}^{B_r \times d}$ and $\mathbf{K}, \mathbf{V} \in \mathbb{R}^{N \times d}$ located in HBM, key block size $B_c$ with $T_c = \lceil \frac{N}{B_c} \rceil$.
      \STATE Adjust the allocation of registers according to the total consumer warps.
      \STATE Set up $\mathbf{O}_i = (0) \in \mathbb{R}^{B_r \times d}$ and initialize $\ell_i, m_i = (0), (-\infty) \in \mathbb{R}^{B_r}$ in local storage.
      \STATE Ensure both $\mathbf{Q}_i$ and $\mathbf{K}_0$ are available in shared memory.
      \STATE Use WGMMA to calculate $\mathbf{S}_{\mathrm{cur}} = \mathbf{Q}_i \mathbf{K}_0^T$. Perform commit and synchronize.
      \STATE Release buffer stage $0$ assigned to $\mathbf{K}$.
      \STATE Derive $m_{i}$, $\tilde{\mathbf{P}}_{\mathrm{cur}}$, and $\ell_{i}$ from $\mathbf{S}_{\mathrm{cur}}$, then update the scaling of $\mathbf{O}_i$.
      \FOR{$1 \le j < T_c - 1$}
        \STATE Synchronize until $\mathbf{K}_j$ is ready in shared memory.
        \label{code:mainloop_start}
        \STATE Calculate $\mathbf{S}_{\mathrm{next}} = \mathbf{Q}_i \mathbf{K}_{j}^T$ via WGMMA. Commit action but remain asynchronous. \label{code:first_wgmma}
        \STATE Wait for $\mathbf{V}_{j-1}$ to arrive in shared memory.
        \STATE Execute $\mathbf{O}_{i} = \mathbf{O}_{i} + \tilde{\mathbf{P}}_{\mathrm{cur}} \mathbf{V}_{j-1}$ using WGMMA, committing but without immediate synchronization. \label{code:second_wgmma}
        \STATE Await completion of the WGMMA calculation for $\mathbf{Q}_i \mathbf{K}_{j}^T$.
        \STATE Update $m_{i}$, obtain $\tilde{\mathbf{P}}_{\mathrm{next}}$, and compute $\ell_{i}$ utilizing $\mathbf{S}_{\mathrm{next}}$. \label{code:softmax}
        \STATE Wait for completion of the WGMMA $\tilde{\mathbf{P}}_{\mathrm{cur}} \mathbf{V}_{j-1}$, followed by adjusting the scaling of $\mathbf{O}_i$.
        \STATE Free the $(j\,\%\,s)$th and $(j-1\,\%\,s)$th buffer stages assigned to $\mathbf{K}$ and $\mathbf{V}$, respectively.
        \STATE Assign $\mathbf{S}_{\mathrm{next}}$ to $\mathbf{S}_{\mathrm{cur}}$.
        \label{code:mainloop_end}
      \ENDFOR
      \STATE Synchronize until $\mathbf{V}_{T_c - 1}$ is available in shared memory.
      \STATE Finalize by calculating 
      $\mathbf{O}_{i} = \mathbf{O}_{i} + \tilde{\mathbf{P}}_{\mathrm{last}} \mathbf{V}_{T_c - 1}$ using WGMMA; commit and synchronize.

\STATE Epilogue: Adjust $\mathbf{O}_{i}$ by applying the scaling factor $m_{i}$. Determine $L_{i}$ utilizing both $m_{i}$ and $\ell_{i}$.  
Store the updated values of $\mathbf{O}_{i}$ and $L_{i}$ in HBM, assigning them as the $i$-th segments of $\mathbf{O}$ and $L$, respectively.  
\end{algorithmic}
\end{algorithm}

\cref{alg:flash3_wgmma} substitutes the consumer portion of \cref{alg:flash3_wgmma_ws_only} to form the full FP16 variant of the \textsc{FlashAttention-3} algorithm. Broadly, the term WGMMA stands in for asynchronous GEMM. During the mainloop (spanning lines \ref{code:mainloop_start} to \ref{code:mainloop_end}), the second WGMMA call in each iteration $j$ (line \ref{code:second_wgmma}) is executed concurrently with the softmax procedures belonging to the subsequent iteration $j+1$ (line \ref{code:softmax}).

Although the pipelined structure outlined above promises theoretical improvements in throughput, there are multiple real-world constraints that must be addressed:

\paragraph{Compiler reordering} The provided pseudocode assumes an ideal order of execution; however, the compiler (NVCC) frequently reorganizes instructions for optimization purposes. This may alter the intended sequence of WGMMA and non-WGMMA operations in the pipeline, potentially causing unexpected execution patterns or reducing anticipated performance benefits. Examination of the generated SASS demonstrates that, in this case, the compiler successfully produces overlapping instructions as intended (see Section~\ref{sec:2-stage-sass}).

\paragraph{Register pressure} Achieving the best performance requires avoiding register spilling. The 2-stage pipeline, however, necessitates more registers to save intermediate data and preserve the context between steps. Specifically, storing an additional $\mathbf{S}_{\mathrm{next}}$ in registers increases the register usage by $B_r \times B_c \times \text{sizeof}(\text{float})$ for each threadblock. This heightened register usage may compete with the need for larger block sizes—which are themselves heavy consumers of registers—as part of other performance optimizations. Empirically, making suitable trade-offs between these requirements depends on profiling.

\paragraph{3-stage pipelining} By building upon the 2-stage procedure, a 3-stage pipelined method is introduced that allows overlapping the second WGMMA computation with the softmax operation. Although this can further enhance Tensor Core utilization, it imposes additional register overhead because of the extra pipeline stage. Balancing tile dimensions against deeper pipeline structures becomes more challenging as a result. A comprehensive discussion and evaluation of the 3-stage algorithm appear in ~\cref{sec:3-stage}.

\subsection{Low-precision with FP8}
\label{sec:algofp8}

\textbf{Efficiency: layout transformations.} Executing the forward pass of \textsc{FlashAttention-3} using FP8 precision introduces unique difficulties related to layout compatibility that are absent when working with FP16.

To begin, it is important to highlight that input tensors $\mathbf{Q}$, $\mathbf{K}$, and $\mathbf{V}$ are commonly arranged with contiguity along the head dimension. However, compliance with the k-major requirement imposed by FP8 WGMMA for the second GEMM necessitates that $\mathbf{V}$—more specifically, the subtiles of $\mathbf{V}$ transferred into SMEM—must instead be contiguous in the sequence length dimension. Given that the TMA loading operation cannot alter which dimension is contiguous, this leads to two possible approaches: (1) performing a transposition of $\mathbf{V}$ within GMEM before computation, or (2) executing an in-kernel tile-wise transpose immediately after moving the data into SMEM. For the first strategy, one can either (1a) integrate the transposition with the epilogue of an earlier process such as rotary embedding, or (1b) invoke a dedicated preprocessing transpose kernel\footnote{A well-optimized transpose kernel can approach the memory bandwidth limit of the hardware~\citep{colfax_cutlass_transpose_2024}.} to swap the strides associated with the sequence length and head axes. Nonetheless, (1a) proves challenging to incorporate into existing library routines, while (1b) is inefficient under conditions where memory bandwidth is the limiting factor, as is the case during inference.

Alternatively, in FP8 \textsc{FlashAttention-3}, we select approach (2). For transposing within the kernel, we utilize the LDSM (\verb|ldmatrix|) and STSM (\verb|stmatrix|) operations. These instructions enable a warp to collaboratively move data between SMEM and RMEM in blocks of 128 bytes. \footnote{The PTX documentation indicates that LDSM/STSM operate by transferring $8 \times 8$ matrices with 16-bit elements \cite[\S 9.7.13.4.15-16]{ptx}, yet, by packing pairs of 8-bit elements, these instructions can be extended to support FP8 computations. Nevertheless, since their transpose variants cannot unpack 8-bit data, additional register-level shuffling is required between LDSM and STSM to achieve the desired tile-wise transpose. We omit these specifics here.} These instructions conserve register resources, which lets us deploy them in the producer warpgroup, and they support layout transposition while copying memory. Furthermore, beginning with the second iteration, we can schedule the transpose of the upcoming $\mathbf{V}$ tile so that it occurs concurrently with the two WGMMA operations that process the previous $\mathbf{V}$ and the current $\mathbf{K}$ tile.

Secondly, we note that for FP8 WGMMA, the FP32 accumulator's memory organization diverges from that of its operand A as represented in the registers—unlike FP16, which maintains a consistent layout. Portions of these memory layouts are illustrated in \cref{fig:rmem_accum} and \cref{fig:rmem_operand}, where register entries are distributed per thread according to the enumerated sequence. Leveraging byte permute instructions, it is possible to reformat the accumulator output of the first WGMMA to match the requirements of a subsequent WGMMA and to align with the structure of the $\mathbf{V}$ tile produced by the in-kernel transpose. More precisely, referencing \cref{fig:rmem_accum}, the ordering is rearranged as
$$\{ \verb|d0 d1 d4 d5 d2 d3 d6 d7| \},$$
and this pattern of register reordering is duplicated across every 8-byte segment. With respect to the conceptual layout of the $\mathbf{P}$ tile, this adjustment permutes its columns, so, for instance, columns $0189$ now occupy the leading four column positions. In order for WGMMA to generate the appropriate output tile, the in-kernel transpose can be implemented such that it produces a corresponding row permutation within the $\mathbf{V}$ tile.\footnote{The flexibility gained by performing an in-kernel transpose removes the necessity for shuffle instructions to transfer register data between threads, as discussed in~\citep{colfax_fp8_flashattention_2024}.}

\textbf{Accuracy: block quantization and incoherent processing.} The FP8 (e4m3) numerical representation allocates just 3 bits for the mantissa and 4 bits for the exponent, thereby incurring greater numerical errors in comparison to FP16 or BF16. Large-scale models often contain outlier values~\citep{dettmers2208llm, sun2024massive} that are significantly more extreme than the majority of entries, which complicates the quantization process. A typical approach is to apply per-tensor scaling~\citep{micikevicius2022fp8}, associating a unique scale factor with each tensor (such as one each for $\mathbf{Q}$, $\mathbf{K}$, and $\mathbf{V}$). To address accuracy degradation in FP8 attention calculations, we adopt two specific strategies:

\begin{enumerate}

\item \textbf{Block quantization}: This approach involves assigning a single scalar value to each block within the tensor. For each of $\mathbf{Q}$, $\mathbf{K}$, and $\mathbf{V}$, the tensors are partitioned into blocks with dimensions $B_r \times d$ or $B_c \times d$, which are then quantized individually. The quantization can be seamlessly integrated with certain operations prior to the attention computation, such as rotary embedding, without incurring any performance penalty (as rotary embedding's speed is restricted by memory bandwidth). Since the \textsc{FlashAttention-3} algorithm inherently works on blocks, it is possible to rescale each block of $\mathbf{S}$ to reflect block quantization, with no extra computational overhead.

\item \textbf{Incoherent processing}: To mitigate the effect of outlier values, both $\mathbf{Q}$ and $\mathbf{K}$ are pre-multiplied with a random orthogonal matrix $\mathbf{M}$ before FP8 quantization. The orthogonality of $\mathbf{M}$ ensures that $\mathbf{M} \mathbf{M}^\top = I$, and thus $(\mathbf{Q} \mathbf{M}) (\mathbf{K} \mathbf{M})^\top = \mathbf{Q} \mathbf{K}^\top$, meaning the attention computation remains unchanged after this transformation. This technique disperses outlier values because every element in $\mathbf{Q} \mathbf{M}$ or $\mathbf{K} \mathbf{M}$ becomes a randomized linear combination of the original entries from $\mathbf{Q}$ or $\mathbf{K}$, thereby diminishing quantization errors. In our experiments, we use the same construction for $\mathbf{M}$ as in \citet{chee2024quip} and \citet{tseng2024quip}, where $\mathbf{M}$ is formed by multiplying random diagonal matrices (with entries $\pm 1$) and a Hadamard matrix, allowing multiplication in $O(d \log d)$ rather than $O(d^2)$. This procedure can also be merged with the rotary embedding step without incurring additional computational cost.
\end{enumerate}

We empirically demonstrate in \cref{sec:numerical_error} that these two strategies achieve up to a 2.6$\times$ reduction in numerical error.

\section{Empirical Validation}
\label{sec:exp}

We leverage CUTLASS~\citep{Thakkar_CUTLASS_2023} primitives, including the WGMMA and TMA abstractions, to construct \textsc{FlashAttention-3} and assess its performance in terms of both efficiency and accuracy.

\begin{itemize}

\item \textbf{Benchmarking attention.} We evaluate the performance of \textsc{FlashAttention-3} by recording its runtime across a range of sequence lengths and contrasting it with a baseline PyTorch implementation, \textsc{FlashAttention-2}, a Triton-based \textsc{FlashAttention-2} leveraging H100-tailored instructions, as well as a cuDNN-optimized vendor implementation for H100 GPUs. Our experiments demonstrate that \textsc{FlashAttention-3} achieves speed improvements of up to 2.0$\times$ compared to \textsc{FlashAttention-2}, and up to 1.5$\times$ relative to the Triton variant. Moreover, \textsc{FlashAttention-3} attains throughput up to 740 TFLOPs/s, amounting to 75\% of the peak theoretical TFLOPs/s on H100 hardware.

\item \textbf{Ablation study.} We provide empirical evidence that enhancements such as warp-specialization and the integration of GEMM-softmax pipelining are responsible for the observed performance gains in \textsc{FlashAttention-3}.

\item \textbf{Accuracy of FP8 attention.} Our analysis demonstrates that employing block quantization together with incoherent processing decreases the numerical error in FP8 \textsc{FlashAttention-3} by a factor of 2.6$\times$.
\end{itemize}

\subsection{Benchmarking Attention}
\label{subsec:benchmark_attn}

We evaluate the performance of various attention algorithms by recording their execution time on an H100 80GB SXM5 GPU, examining different configurations (with and without causal masking, head dimensions set to 64 or 128) using FP16 input tensors. Our findings, presented in~\cref{fig:benchmark_attn_fwd} and~\cref{fig:benchmark_attn_bwd}, indicate that \textsc{FlashAttention-3} delivers a speedup of approximately 1.5-2.0$\times$ over \textsc{FlashAttention-2} in the forward pass and achieves about 1.5-1.75$\times$ acceleration in the backward pass. Relative to a conventional attention baseline, \textsc{FlashAttention-3} attains speedups between 3$\times$ and 16$\times$. Furthermore, for sequences of moderate to considerable length (at least 1k tokens), \textsc{FlashAttention-3} outperforms even the proprietary vendor-optimized cuDNN library designed specifically for the H100 GPUs.

\paragraph{Benchmark settings:} We vary the sequence length as 512, 1k, ..., 16k, and set batch size so that the total number of tokens is 16k. We set the hidden dimension to 2048, and head dimension to be either 64, 128, or 256 (i.e., 32 heads, 16 heads, or 8 heads). To calculate the FLOPs of the forward pass, we use:

\begin{equation*}
  4 \cdot \text{seqlen}^2 \cdot \text{head dimension} \cdot \text{number of heads}.
\end{equation*}

By employing causal masking, this value is halved to reflect that roughly 50% of the elements are computed. For the backward pass FLOPs, we scale the forward pass FLOPs by a factor of 2.5. This stems from the fact that while the forward computation involves 2 matrix multiplications, the backward operation entails 5 matrix multiplications, owing to recomputation.

Additionally, we assess the forward pass runtime using FP8 under comparable conditions. The outcomes for headdim 256 are displayed in ~\cref{fig:benchmark_attn_fp8}, while comprehensive findings can be found in \cref{sec:benchmark_attn_fp8_full}.

\subsection{Ablation Study: 2-Stage Pipelining Experiments}
\label{sec:2-stage-pipelining-experiments}

We conduct ablation studies on both the two-stage WGMMA-softmax pipelining and the warp-specialization techniques within the non-causal FP16 \textsc{FlashAttention-3} framework, while holding parameters constant at $\{\text{batch}, \text{seqlen}, \text{nheads}, \text{hdim}\} = \{ 4, 8448, 16, 128\}$. The findings, as presented in~\cref{table:ablation_pipelining}, demonstrate that incorporating algorithmic enhancements—namely, asynchronous execution via warp-specialization and concurrent processing of GEMM and softmax—yields a substantial performance boost, raising throughput from 570 to 661 TFLOPs.

\subsection{Numerical Error Validation}
\label{sec:numerical_error}

Given ongoing concerns regarding numerical errors in \textsc{FlashAttention}~\citep{golden2024flash}, we evaluate \textsc{FlashAttention-2}, \textsc{FlashAttention-3}, as well as a conventional attention implementation by benchmarking them against a FP64 reference implementation. To replicate the presence of outlier features and activations observed in LLMs~\citep{dettmers2208llm, sun2024massive}, we draw the elements of $\mathbf{Q}, \mathbf{K}, \mathbf{V}$ from the distribution described below:

\begin{equation*}
  \mathcal{N}(0, 1) + \mathcal{N}(0, 100) \cdot \mathrm{Bernoulli}(0.001).
\end{equation*}

Specifically, each element is sampled from a normal distribution with mean zero and a standard deviation of 1, except for 0.1\% of the elements, to which we add an independent Gaussian noise term with a standard deviation of 10. The root mean squared error (RMSE) for this setup is reported in~\cref{table:numerical_error}. In FP16, both \textsc{FlashAttention-2} and \textsc{FlashAttention-3} attain an RMSE that is 1.7$\times$ smaller than the conventional approach due to the retention of intermediate (softmax) calculations in FP32 precision. For the FP8 baseline, per-tensor scaling is applied, with matrix multiplications accumulated in FP32, and intermediate softmax outputs preserved in FP16. Owing to block quantization and the use of incoherent processing, \textsc{FlashAttention-3} in FP8 demonstrates a 2.6$\times$ improvement in accuracy over the aforementioned baseline.

\section{Dicussion, Limitations, Conclusion}
\label{sec:discussion}

Utilizing \textsc{FlashAttention-3}, our findings illustrate that adopting novel programming strategies alongside modern hardware features like asynchrony and reduced precision significantly improves both the efficiency and fidelity of attention mechanisms. Relative to \textsc{FlashAttention-2}, attention computation achieves a speedup ranging from 1.5$\times$ to 2.0$\times$, while FP8 numerical error is lowered by a factor of 2.6$\times$ compared to conventional per-tensor quantization approaches. There remain several limitations in the current study that warrant further investigation: enhancing LLM inference performance, implementing a persistent kernel architecture within the FP8 kernel,\footnote{Our benchmarking setup includes a persistent kernel and load balancing for FP16 \textsc{FlashAttention-3}, whereas FP8 \textsc{FlashAttention-3} lacks these optimizations. This disparity partly accounts for its inferior performance under small sequence lengths and causal masking, relative to FP8 cuDNN kernels.} and evaluating the ramifications of low-precision attention during large-scale training. While this research centers on Hopper GPUs, it is anticipated that these methodologies are generalizable to alternative hardware accelerators. The expectation is that improved, faster attention primitives may catalyze progress in applications requiring extended context.

\end{document}